% Resume: Arumalla Mohit Krishna
% Compile with pdfLaTeX.

\documentclass[letterpaper,10pt]{article}

\usepackage[margin=0.5in]{geometry}
\usepackage[T1]{fontenc}
\usepackage{lmodern}
\usepackage{microtype}
\usepackage{titlesec}
\usepackage{enumitem}
\usepackage[hidelinks]{hyperref}
\usepackage{tabularx}
\input{glyphtounicode}

\pdfgentounicode=1
\pagestyle{empty}
\urlstyle{same}
\raggedright
\setlength{\parindent}{0pt}
\setlength{\tabcolsep}{0pt}

\titleformat{\section}
  {\large\scshape}{}{0em}{}[\titlerule]

\titlespacing*{\section}{0pt}{6pt}{3pt}

\newcommand{\entry}[2]{%
  \noindent
  \begin{tabularx}{\linewidth}
    {@{}>{\raggedright\arraybackslash}X@{\hspace{6pt}}r@{}}
    \textbf{#1} & #2
  \end{tabularx}\par
}

\newcommand{\bullets}{
  \begin{itemize}[
    leftmargin=1.2em,
    topsep=2pt,
    itemsep=1pt,
    parsep=0pt,
    partopsep=0pt
  ]
}

\newcommand{\finishbullets}{
  \end{itemize}
}

\begin{document}

%----------HEADING----------
\begin{center}
    {\LARGE\bfseries Arumalla Mohit Krishna}\\[3pt]
    +1 (631) 352-8643 $|$
    \href{mailto:mohitkrishna.arumalla@stonybrook.edu}
    {mohitkrishna.arumalla@stonybrook.edu} $|$
    \href{https://linkedin.com/in/mohit-krishna-arumalla}
    {LinkedIn}
\end{center}
\vspace{-8pt}

%-----------EDUCATION-----------
\section{Education}

\entry
  {\textbf{Stony Brook University}}
  {Aug 2025 -- Dec 2026}

M.S. in Computer Science \hfill New York, USA\\
\textit{
Coursework: \textbf{Natural Language Processing},
Introduction to Computer Vision, Vision Language Modeling
}
\vspace{3pt}

\entry
  {\textbf{Birla Institute of Technology and Science, Pilani}}
  {Aug 2019 -- 2023}

B.E. in Electronics \& Communication; GPA: 8.02/10 \hfill India
%-----------PROFESSIONAL SUMMARY-----------
\section{Professional Summary}

Software engineer with 2+ years of experience in \textbf{AI and machine learning solution development},
building distributed services and real-time data pipelines. Currently pursuing an M.S.
in Computer Science at Stony Brook University with specialized knowledge in \textbf{Generative AI},
deep learning frameworks, and scalable \textbf{data processing pipelines}. Skilled in Python, PyTorch,
TensorFlow, Kubernetes, and cloud infrastructure with strong collaboration and communication abilities.
%-----------TECHNICAL SKILLS-----------
\section{Technical Skills}

\textbf{Languages:}
Python, Go, C++, SQL, Java\\

\textbf{ML \& AI Frameworks:}
PyTorch, TensorFlow/Keras, scikit-learn, NumPy, pandas,
Transformers, Generative AI, Deep Learning,
Self-Supervised Learning, Vision Transformers,
LoRA, Visual Prompt Tuning, Text Embeddings, Vector Search\\

\textbf{Data Engineering \& Systems:}
gRPC, Kafka, Apache Flink, Flink SQL, Debezium/CDC,
PostgreSQL, Redis, Elasticsearch, Airflow, BigQuery,
\textbf{Kubernetes}, HPC clusters, Spark\\

\textbf{Infrastructure:}
NVIDIA Triton, AWS, Docker, Kubernetes, Linux, Git,
CI/CD, Istio/Envoy, New Relic, Cloud Native Tools

%-----------EXPERIENCE-----------
\section{Experience}

\entry
  {\textbf{Software Engineer II} $|$ \textbf{Infra.Market, Bangalore}}
  {Apr 2024 -- Aug 2025}

\bullets

  \item Designed and implemented scalable AI/ML applications integrating Amazon Titan text embeddings
  with Elasticsearch vector search for enhanced semantic product discovery.

  \item Engineered a high-throughput CDC pipeline using Flink, Kafka, Debezium,
  PostgreSQL, and Elasticsearch, enabling near real-time data synchronization (<500 ms latency)
  with strong data consistency guarantees.

  \item Deployed and maintained Kubernetes clusters with Istio/Envoy service mesh,
  integrating distributed tracing via New Relic for secure and robust production monitoring.

\finishbullets

\entry
  {\textbf{Software Engineer I} $|$ \textbf{Infra.Market, Chennai}}
  {Jul 2023 -- Mar 2024}

\bullets

  \item Owned development of Search and Location microservices; created performant gRPC services
  handling 500K+ daily requests, supporting critical business workflows and geospatial features.

  \item Resolved PostgreSQL connection pool leaks using Gorm ORM, improving system stability
  and throughput by 85\%; validated fixes with high-concurrency load testing.

  \item Automated Airflow pipelines for data migration from BigQuery to Elasticsearch with
  dynamic filtering and stale-index monitoring to maintain data quality and freshness.

\finishbullets

%-----------PROJECTS-----------
\section{Projects}

\entry
  {GPU-Optimized Personalized Product Search}
  {Aug 2026 -- Oct 2026}

\textit{Python, PyTorch, BGE, NVIDIA Triton, Go, gRPC}

\bullets

  \item Developed two-tower retrieval model using \texttt{BAAI/bge-small-en-v1.5} embeddings for
  query and product, combining with user interaction histories to personalize search.

  \item Built scalable inference pipeline with precomputed embeddings, approximate nearest-neighbor
  search, and GPU-accelerated query encoding via NVIDIA Triton, supporting low-latency production use.

  \item Benchmarked throughput, latency, and resource utilization under varied traffic,
  optimizing for \textbf{scalable and robust AI inference} serving with recall and NDCG metrics.

\finishbullets

\entry
  {Attention-Based Hallucination Detection in NMT}
  {Apr 2026 -- May 2026}

\textit{Python, PyTorch, mBART, Transformers, scikit-learn}

\bullets

  \item Built English-German hallucination detection pipeline leveraging teacher-forced mBART attention
  features and WMT22/WMT23 quality estimation datasets.

  \item Extracted detailed cross/self-attention metrics across decoder layers and heads for deep model interpretability.

  \item Trained classifiers achieving 0.712 F1 and 0.695 ROC-AUC; proposed hallucination-aware reranking
  methods integrating attention-based risk scores.

\finishbullets

\entry
  {Prompt/LoRA for Self-Supervised Learning}
  {Ongoing}

\textit{Python, PyTorch, Vision Transformers}

\bullets

  \item Researching domain-specialized visual prompt tuning and parameter-efficient LoRA adapters within
  self-supervised ViT architectures to improve cross-domain transferability.

  \item Evaluating performance impacts of modular adaptation methods on classification accuracy
  and parameter efficiency under domain shifts.

\finishbullets

\end{document}